%% kapa: exam=1 round=32 year=2021 period=2 form=A subject=accounting subject_ko=회계학 seq=21 number=61
%% source: Q-Net 감정평가사 기출문제 (cst00309) artlSeq=5208397 file=extracted/1008/5208397/A형/2021년도 제32회 감정평가사 1차 2교시 A형.pdf
%% conversion: kordoc 4.24.51 -> tools/gen1.py (auto); review=proofread-sonnet; flags=none
\begin{question}{61}
고객과의 계약으로 식별하기 위한 기준에 관한 설명으로 옳지 않은 것은?
\begin{choices}
  \item 계약 당사자들이 계약을 서면으로, 구두로 또는 그 밖의 사업 관행에 따라 승인하고 각자의 의무를 수행하기로 확약한다.
  \item 이전할 재화나 용역과 관련된 각 당사자의 권리를 식별할 수 있다.
  \item 이전할 재화나 용역의 지급조건을 식별할 수 있다.
  \item 계약에 상업적 실질을 요하지는 않는다.
  \item 고객에게 이전할 재화나 용역에 대하여 받을 권리를 갖게 될 대가의 회수 가능성이 높다.
\end{choices}
\end{question}
\answer{4}
